
%%%%%%%%%%%%%%%%%%%%%%%%%%%%%%
%% 1. Table and Figure Settings
%%%%%%%%%%%%%%%%%%%%%%%%%%%%%%
% Format subfigure captions.
% Result: Separates label and caption with a space; sets 1pt vertical skip.
% Note: If "Unused \captionsetup[subfigure]" warning occurs, subcaption package might be missing or conflicting. Keep commented if unused.
%\captionsetup[subfigure]{labelsep=space, skip=1pt}

%%%%%%%%%%%%%%%%%%%%%%%%%%%%%%
%% 2. Algorithm Customizations
%%%%%%%%%%%%%%%%%%%%%%%%%%%%%%
% Note: The following settings rely on 'algorithm2e'.
% Since 'icml2026.sty' loads conflicting packages by default, we have commented out algorithm2e in packages.tex.
% Therefore, these settings must also be commented out to avoid errors.

% Define 'Input' keyword
% Result: Displays "Input: " in algorithms.
% \SetKwInput{KwInput}{Input}

% Define 'Output' keyword
% Result: Displays "Output: " in algorithms.
% \SetKwInput{KwOutput}{Output}

% Define 'Block' comment style
% Result: Uses '⊳' symbol for inline comments.
% \SetKwComment{Block}{$\triangleright$\ }{}

% Attempt to fix compatibility between hyperref and algorithmic
% \newcommand{\theHalgorithm}{\arabic{algorithm}}

% Define standard comment style
% Result: Uses '//' for comments.
% \SetKwComment{Comment}{//  }{}

% Prevent automatic semicolon printing
% Result: No semicolons automatically added at line ends.
% \DontPrintSemicolon

% Customize 'if-else' structure
% Result: Sets display keywords for If, ElseIf, Else.
% \SetKwIF{If}{ElseIf}{Else}{if}{}{else if}{else}{end if}%

%%%%%%%%%%%%%%%%%%%%%%%%%%%%%%
%% 3. Highlighting and Colors
%%%%%%%%%%%%%%%%%%%%%%%%%%%%%%
% Define line highlighting command \HiLi
% Result: Creates a light gray background for the current line (useful for emphasizing table rows).
\def\HiLi{\leavevmode\rlap{\hbox to \hsize{\color{gray!20}\leaders\hrule height .8\baselineskip depth .5ex\hfill}}}

% Define dark blue color
\definecolor{darkblue}{rgb}{0, 0, 0.5}

% Configure hyperlink colors
% Result: Citations, internal links, and URLs appear in dark blue and are clickable in PDF.
\hypersetup{
  colorlinks=true,   % Enable colored links
  citecolor=darkblue, % Color for citations (e.g., [1])
  linkcolor=darkblue, % Color for internal links (e.g., Figure 1)
  urlcolor=darkblue   % Color for external URLs
}

%%%%%%%%%%%%%%%%%%%%%%%%%%%%%%
%% 4. Code Listings Configuration
%%%%%%%%%%%%%%%%%%%%%%%%%%%%%%
% Configure 'listings' package display
\lstset{
  basicstyle=\ttfamily\footnotesize,       % Font: Monospaced, footnote size
  breaklines=true,                         % Line breaking: Auto-wrap long lines
  breakatwhitespace=false,                 % Break at: Allow breaks at any character, not just spaces
  showstringspaces=false,                  % String spaces: Do not show visible markers for spaces in strings
  columns=flexible,                        % Alignment: Flexible character width to avoid awkward spacing
  frame=single,                            % Frame: Add a single-line frame around code blocks
  backgroundcolor=\color{lightgray!20}     % Background: Light gray background for code blocks
}

%%%%%%%%%%%%%%%%%%%%%%%%%%%%%%
%% 5. Theorem and Proof Environments
%%%%%%%%%%%%%%%%%%%%%%%%%%%%%%
% Set theorem style to 'plain'
% Result: Bold title, italic body (for core statements like Theorems, Lemmas)
\theoremstyle{plain}

% Define 'theorem' environment
% Result: Numbered by section (e.g., Theorem 1.1)
\newtheorem{theorem}{Theorem}[section]

% Define 'proposition' environment
% Result: Shares numbering with Theorem (e.g., Proposition 1.2)
\newtheorem{proposition}[theorem]{Proposition}

% Define 'lemma' environment
% Result: Shares numbering with Theorem (e.g., Lemma 1.3)
\newtheorem{lemma}[theorem]{Lemma}

% Define 'corollary' environment
% Result: Shares numbering with Theorem (e.g., Corollary 1.4)
\newtheorem{corollary}[theorem]{Corollary}

% Set theorem style to 'definition'
% Result: Bold title, upright body (for Definitions, Assumptions)
\theoremstyle{definition}

% Define 'definition' environment
% Result: Shares numbering with Theorem
\newtheorem{definition}[theorem]{Definition}

% Define 'assumption' environment
% Result: Shares numbering with Theorem
\newtheorem{assumption}[theorem]{Assumption}

% Set theorem style to 'remark'
% Result: Italic title, upright body (for Remarks)
\theoremstyle{remark}

% Define 'remark' environment
% Result: Shares numbering with Theorem
\newtheorem{remark}[theorem]{Remark}



